% ============================================================
% Титульный лист
% ============================================================

% Встроенный титульный лист
\titlepagestyle{system/titlepages/minimal}

% Свой сверстанный титульный лист
% \titlepagestyle{user/titlepages/titlepage}

% Внешний титульный лист PDF
% \titlepagepdf{assets/titlepage.pdf}


% ============================================================
% Содержание
% ============================================================

\ConfigShowTableOfContentstrue
% \ConfigShowTableOfContentsfalse


% ============================================================
% PDF
% ============================================================

% Полный список доступных параметров:
% https://mirrors.ctan.org/macros/latex/contrib/hyperref/doc/hyperref-doc.pdf

\hypersetup{
    hidelinks,
    bookmarksnumbered=true,
    breaklinks=true,
    pdfstartview=Fit
}

% ------------------------------------------------------------
% Разделы
% ------------------------------------------------------------

% Каждый \section начинается с нового листа.
\ConfigSectionNewPagetrue
% \ConfigSectionNewPagefalse

% ============================================================
% Библиография
% ============================================================

% Показывать список источников.
\ConfigShowBibliographytrue
% \ConfigShowBibliographyfalse

% Нумеровать список литературы.
\ConfigBibliographyNumberedtrue
% \ConfigBibliographyNumberedfalse

% Включать все источники в список литературы,
% даже если они не цитируются.
\ConfigBibliographyAlltrue
% \ConfigBibliographyAllfalse

% Файл библиографии.
\renewcommand{\ConfigBibliographyFile}{src/bibliography/sources.bib}

\addbibresource{\ConfigBibliographyFile}

% ============================================================
% Листинг
% ============================================================

% ------------------------------------------------------------
% Нумерация листингов
% ------------------------------------------------------------

% Сквозная:
% Листинг 1, Листинг 2, Листинг 3, ...
%\ConfigListingNumberedBySectionfalse

% По разделам:
% Листинг 1.1, Листинг 1.2, Листинг 2.1, ...
 \ConfigListingNumberedBySectiontrue

% ------------------------------------------------------------
% Название листинга
% ------------------------------------------------------------

% Подпись:
% «Листинг 1 — Пример программы»
% «Листинг 1.1 — Пример программы»
\renewcommand{\lstlistingname}{Листинг}

% ------------------------------------------------------------
% Подписи листингов
% ------------------------------------------------------------

% Размер шрифта подписи.
\newcommand{\ConfigListingCaptionFont}{small}

% Выравнивание подписи.
% centering, raggedright, raggedleft
\newcommand{\ConfigListingCaptionJustification}{centering}

% Разделитель между номером и названием:
% endash — «Листинг 1 — Пример»
% space  — «Листинг 1 Пример»
% colon  — «Листинг 1: Пример»
\newcommand{\ConfigListingCaptionLabelSep}{colon}

% Центрирование коротких подписей.
\newcommand{\ConfigListingCaptionSingleLineCheck}{true}

% Применить настройки.
\ESKDconfigureListingCaption


% ============================================================
% Рисунки
% ============================================================

% ------------------------------------------------------------
% Каталог изображений
% ------------------------------------------------------------

\graphicspath{{assets/images/}}

% ------------------------------------------------------------
% Размер изображений
% ------------------------------------------------------------

% Максимальная ширина изображения.
\newcommand{\ConfigFigureWidth}{0.8\textwidth}

% Максимальная высота изображения.
\newcommand{\ConfigFigureHeight}{0.3\textheight}

% ------------------------------------------------------------
% Подписи рисунков
% ------------------------------------------------------------

% Размер шрифта подписи.
% tiny, scriptsize, footnotesize, small,
% normalsize, large, Large, LARGE, huge, Huge
\newcommand{\ConfigFigureCaptionFont}{small}

% Выравнивание подписи.
% centering, raggedright, raggedleft
\newcommand{\ConfigFigureCaptionJustification}{centering}

% Разделитель между номером и названием:
% endash — «Рис. 1.1 — Название»
% space  — «Рис. 1.1 Название»
% colon  — «Рис. 1.1: Название»
\newcommand{\ConfigFigureCaptionLabelSep}{endash}

% Центрирование коротких подписей.
\newcommand{\ConfigFigureCaptionSingleLineCheck}{true}

% ------------------------------------------------------------
% Нумерация рисунков
% ------------------------------------------------------------

% Сквозная нумерация:
% Рис. 1, Рис. 2, Рис. 3, ...

%\ConfigFigureNumberedBySectionfalse

% Нумерация по разделам:
% Рис. 1.1, Рис. 1.2, Рис. 2.1, ...

 \ConfigFigureNumberedBySectiontrue

\ifConfigFigureNumberedBySection
\counterwithin{figure}{section}
\fi

% ------------------------------------------------------------
% Настройки рамки
% ------------------------------------------------------------

\newif\ifConfigFigureFrameEnabled
\ConfigFigureFrameEnabledfalse

% ------------------------------------------------------------
% Рамка изображения
% ------------------------------------------------------------

% Показывать рамку вокруг изображения.
\ConfigFigureFrameEnabledtrue
% \ConfigFigureFrameEnabledfalse

% Цвет рамки.
\definecolor{ConfigFigureFrameColor}{RGB}{180,180,180}

% Толщина рамки.
\newcommand{\ConfigFigureFrameRule}{0.8pt}

% Отступ между изображением и рамкой.
\newcommand{\ConfigFigureFrameSep}{0pt}

% Применить настройки.
\ESKDconfigureFigureCaption